% ============================================================================
% 00_abstract.tex — merged paper (auction-v2)
% \begin{abstract}...\end{abstract} only; no citations (per CONVENTIONS §1).
% Auctions = main domain; matching demoted to one appendix clause.
% Menu claim per F1; clock-framing cross-model per F2; Section 7 (ordinal vs
% cardinal) per F9; traces dissociation per F8.
% ============================================================================
\begin{abstract}
Incentive compatibility on paper does not guarantee dominant-strategy play in practice: agents must recognize that truthful play is safe. As large language model (LLM) agents begin to transact, a design question becomes operational: which levers---the extensive form, the mechanism's description, or scaffolds for the agent's reasoning---preserve truthful play for this new bounded reasoner? We rank these levers empirically in auction environments across four LLM families, using a testbed validated against moment-matched reconstructions of human experimental data: LLM bidders reproduce the human difficulty ranking of auction formats, winner's-curse comparative statics, and last-minute bidding under hard-close rules. The levers form a stable tier structure. Obviously strategy-proof extensive forms rank at the top. Safety-exposing descriptions, clock-framing---which improves play in all four models---and contingent-reasoning scaffolds form a strong second tier; menu restatements have no consistent effect---small, sign-mixed across models, null on average; forward-planning and belief-formation scaffolds actively backfire. The active ingredient is a description's invariance content: text exposing why truth-telling is safe helps, while mechanical restatement does not---and reasoning traces show the safety description moves bids without changing stated reasoning. Human evidence agrees in sign wherever a benchmark exists, and the results replicate in a school-choice matching domain (appendix). Yet the agreement is ordinal, not cardinal: no single tuning of the agents makes their bidding match human behavior in levels across mechanisms.
\end{abstract}
